\section{Conclusion}
\label{sec:conclusion}

This paper develops a rigorous copula-based theory of mean-reverting Markov
processes and examines its empirical relevance for USDT/USD.  The theory makes
two distinctions that are essential in applications.  First, an invariant
distribution or a decaying dependence coefficient is not, by itself, a
directional mean-reversion result; the conditional drift must point toward the
target.  Second, a finite-dimensional formula is not an observable model unless
its transformation is invertible on the relevant support and its transition
density is normalized.  These requirements preserve the valid affine and
monotone OU/CIR constructions while explaining the empirical failure of several
quadratic, exponential, and special-function branches.

The mixed-copula construction is the main theoretical result.  A multiplicative
mixture of a stationary transition copula and independence defines a consistent
Markov semigroup, leaves the invariant marginal unchanged, and has a Poisson
reset representation.  Its conditional centered mean is the base conditional
mean multiplied by the no-reset probability, so the local generator acquires
the transparent restoring term $\lambda(m-x)$.  This provides a parsimonious
way to alter transition density shape and recovery speed without rebuilding the
marginal distribution.

The stablecoin evidence is consistent with the motivation for that separation,
but it also defines its limits.  USDT/USD exhibits finite-horizon directional
reversion, while a single OU transition fails density diagnostics.  The mixed-OU
kernel is strongly favored in the matched in-sample boundary experiment, and
reset mixtures improve out-of-sample log scores for Gaussian and two-scale
Gaussian dependence.  The improvement is not generic across copula families,
and a constant reset rate cannot reproduce the observed slow multi-horizon
persistence.  The strongest five-minute forecasting model is instead a
two-state OU with weakly identified slow-regime dynamics.  Accordingly, the
evidence supports the mixed copula as a structurally coherent density extension,
not as a universally dominant forecasting law or a uniquely identified model of
redemption and arbitrage.

Several extensions follow naturally.  Economic interpretation of reset risk
requires covariates for liquidity, redemptions, reserve news, and market stress.
Slow persistence calls for an explicitly augmented Markov state rather than a
horizon-by-horizon scalar copula approximation.  Finally, the nonlocal pricing
equation should be taken to bond or derivative data under a specified
risk-neutral measure; it is not identified by a spot stablecoin series alone.
These extensions preserve the paper's central principle: marginal behavior,
temporal dependence, mean-reversion direction, and predictive performance must
be modeled and tested separately.
